\begin{enumerate}[label=\thesubsection.\arabic*, ref=\thesubsection.\theenumi]
	
% =========================================================
% Q.3
% [Reference: General Domain Knowledge - Matrix indexing / Indicator variables]
% =========================================================
\item Indicator Summation for Thresholding
\begin{align}
	N = \sum_{i=1}^M \sum_{j=1}^N \mathbb{I}\brak{U_{ij}} \label{eq:DA2025-threshold}
\end{align}
\solution
The total count of elements in a matrix $\vec{U}$ satisfying a threshold $c$ is obtained by summing the indicator function $\mathbb{I}$. The values of $\mathbb{I}$ operate as a step function defined as:
\begin{align}
	\mathbb{I}\brak{U_{ij}} = \begin{cases} 
		1, & U_{ij} \le c \\ 
		0, & U_{ij} > c 
	\end{cases} \label{eq:DA2025-indicator_step}
\end{align}
	
	% =========================================================
	% Q.5
	% [Reference: General Domain Knowledge - Basic Geometry]
	% =========================================================
	\item Geometric Scaling under Dilation
	\begin{align}
		A' = \brak{1 + \alpha}A \label{eq:DA2025-scaling}
	\end{align}
	\solution
	For a rectangle with length $L$ and width $W$, Area $A = LW$. If $W$ is scaled by a factor of $\brak{1 + \alpha}$:
	\begin{align}
		W' &= \brak{1 + \alpha}W \label{eq:DA2025-width_scaled}
	\end{align}
	Substituting \eqref{eq:DA2025-width_scaled} into the area equation:
	\begin{align}
		A' &= L W' = L \sbrak{\brak{1 + \alpha}W} = \brak{1 + \alpha}LW = \brak{1 + \alpha}A
	\end{align}
	yielding \eqref{eq:DA2025-scaling}.
	
	% =========================================================
	% Q.7
	% [Reference: bt-2025-formulae.tex (Item 6: Convexity & Minima definitions)]
	% =========================================================
	\item Injectivity of Strictly Unimodal Functions
	\begin{align}
		x_1 < c < x_2 \quad \text{for} \quad f\brak{x_1} = f\brak{x_2} \label{eq:DA2025-unimodal}
	\end{align}
	\solution
	If $f\brak{x}$ strictly increases on $[0, c]$ and strictly decreases on $[c, \infty)$, each branch is individually injective (one-to-one). If two distinct points $x_1$ and $x_2$ yield the same function value, they cannot lie on the same monotonic branch, meaning the peak $c$ must lie strictly between them.
	
	% =========================================================
	% Q.9
	% [Reference: formulae.tex (Item 14: Complex Logarithm / Exponents)]
	% =========================================================
	\item Exponent Laws
	\begin{align}
		\frac{a^{u}}{a^{v}} = a^{u-v}, \quad \brak{a^x}^y = a^{xy} \label{eq:DA2025-exponent}
	\end{align}
	
	% =========================================================
	% Q.11
	% [Reference: ce2-2025-problems.tex (Q.51: Continuous Expected Values)]
	% =========================================================
	\item Law of Total Expectation (Tower Property)
	Let $X$ and $Y$ be continuous random variables. The initial integral equations defining the expectations are:
	\begin{align}
		g\brak{y} = E\sbrak{X \mid Y = y} &= \int_{-\infty}^{\infty} x f_{X \mid Y}\brak{x \mid y} \, dx \label{eq:DA2025-cond_exp} \\
		E\sbrak{E\sbrak{X \mid Y}} &= \int_{-\infty}^{\infty} g\brak{y} f_Y\brak{y} \, dy \label{eq:DA2025-tower_integral} \\
		E\sbrak{X} &= \int_{-\infty}^{\infty} x f_X\brak{x} \, dx \label{eq:DA2025-expected_x}
	\end{align}
	% =========================================================
	% Q.12
	% [Reference: bt-2025-formulae.tex (Item 8: System of Linear Eq / Matrices)]
	% =========================================================
	\item Computational Complexity of Back-Substitution
	For an $n \times n$ upper triangular matrix system $\vec{U}\vec{x} = \vec{b}$, variables are obtained sequentially:
	\begin{align}
		x_i &= \frac{1}{U_{ii}}\brak{b_i - \sum_{j=i+1}^n U_{ij} x_j}
	\end{align}
	
% =========================================================
% Q.13
% [Reference: bt-2025-formulae.tex (Item 8: System of Linear Eq & Rouché-Capelli Theorem)]
% =========================================================
\item Eigenvalues of Constant Row-Sum Matrices and System Solvability
\begin{align}
	\vec{A}\vec{1} = 1 \cdot \vec{1}, \quad \brak{\vec{A} - \vec{I}}\vec{1} = 0 \cdot \vec{1} \label{eq:DA2025-row_sum_eig}
\end{align}
\solution
Let $\vec{a}_i^\top$ denote the $i$-th row of $\vec{A} \in \mathbb{R}^{n \times n}$. Because the sum of the elements in each row is $1$, multiplying the row by the all-ones column vector $\vec{1} = \myvec{1 & 1 & \dots & 1}^\top$ yields:
\begin{align}
	\vec{a}_i^\top \vec{1} = 1, \quad \forall i \in \{1, 2, \dots, n\}
\end{align}
Grouping these row equations together to reconstruct the full matrix $\vec{A}$ gives:
\begin{align}
	\vec{A}\vec{1} &= \myvec{\vec{a}_1^\top \\ \vec{a}_2^\top \\ \vdots \\ \vec{a}_n^\top} \vec{1} 
	= \myvec{\vec{a}_1^\top \vec{1} \\ \vec{a}_2^\top \vec{1} \\ \vdots \\ \vec{a}_n^\top \vec{1}} 
	= \myvec{1 \\ 1 \\ \vdots \\ 1} = \vec{1} = 1 \cdot \vec{1}
\end{align}
This demonstrates that the vector $\vec{1}$ is an eigenvector of $\vec{A}$ associated with the eigenvalue $\lambda = 1$. 

To find the eigenvalue of $\vec{A} - \vec{I}$ corresponding to the same eigenvector $\vec{1}$:
\begin{align}
	\brak{\vec{A} - \vec{I}}\vec{1} = \vec{A}\vec{1} - \vec{I}\vec{1} = \vec{1} - \vec{1} = \vec{0} = 0 \cdot \vec{1}
\end{align}
Thus, $\vec{A} - \vec{I}$ possesses an eigenvalue of $0$, meaning it is a singular matrix ($\mydet{\vec{A} - \vec{I}} = 0$).

Number of Solutions for Linear Systems:

\begin{enumerate}
	\item For the system $\brak{\vec{A} - \vec{I}}\vec{x} = \vec{b}$:
	Because $\vec{A} - \vec{I}$ is singular ($\text{rank} < n$), the system will have infinitely many solutions if $\vec{b}$ is in the column space of $\vec{A} - \vec{I}$ (i.e., $\text{rank}\brak{\vec{A} - \vec{I}} = \text{rank}\brak{\sbrak{\vec{A} - \vec{I} \mid \vec{b}}}$), or no solution if it is not.
	
	\item For the system $\vec{A}\vec{x} = \vec{b}$:
	The number of solutions depends entirely on the remaining unknown eigenvalues of $\vec{A}$. If all other eigenvalues are non-zero ($\mydet{\vec{A}} \neq 0$), the system has a unique solution. If $0$ happens to also be an eigenvalue of $\vec{A}$, the matrix is singular, resulting in either no solution or infinitely many solutions
\end{enumerate}
	% =========================================================
	% Q.14
	% [Reference: bt-2025-formulae.tex (Item 1: Taylor and Maclaurin Series Foundations)]
	% =========================================================
	\item Maclaurin Series of Hyperbolic Sine
	\begin{align}
		\sinh\brak{x} = \sum_{k=0}^{\infty} \frac{x^{2k+1}}{\brak{2k+1}!} \label{eq:DA2025-maclaurin}
	\end{align}
	\solution
	Expanding $e^x$ and $e^{-x}$ via Maclaurin series :
	\begin{align}
		f\brak{x} &= \frac{1}{2}\brak{\sum_{n=0}^{\infty} \frac{x^n}{n!} - \sum_{n=0}^{\infty} \frac{\brak{-1}^n x^n}{n!}}
		\\
		&= x + \frac{x^3}{3!} + \frac{x^5}{5!} + \cdots
	\end{align}
	
% =========================================================
% Q.17
% [Reference: General Domain Knowledge - SQL/Relational Algebra]
% =========================================================
\item Relational Algebra Natural Join
% natural_join.tex
\begin{figure}[H]
	\centering
	\begin{tikzpicture}[
		table/.style={
			matrix of nodes,
			nodes={draw, minimum width=1.5cm, minimum height=0.6cm, align=center},
			column sep=-\pgflinewidth,
			row sep=-\pgflinewidth,
			nodes in empty cells
		},
		header/.style={fill=gray!20},
		>=Stealth
		]
		
		% Table R
		\matrix[table, label=above:{Relation $R$}] (R) {
			|[header]| A & |[header]| B \\
			1 & X \\
			2 & Y \\
			2 & U \\
			4 & V \\
		};
		
		% Table S
		\matrix[table, right=2.5cm of R, label=above:{Relation $S$}] (S) {
			|[header]| A & |[header]| C \\
			1 & Z \\
			2 & W \\
			3 & T \\
		};
		
		% Join Symbol
		\path (R.east) -- (S.west) node[midway] (join) {\Large $\bowtie$};
		
		% Result Table
		\matrix[table, below=1.5cm of join, label=above:{Result: $R \bowtie S$}] (Result) {
			|[header]| A & |[header]| B & |[header]| C \\
			1 & X & Z \\
			2 & Y & W \\
			2 & U & W \\
		};
		
		% Arrows tracking the match
		\draw[->, thick, dashed] (R.south) -- (Result.north west);
		\draw[->, thick, dashed] (S.south) -- (Result.north east);
		
		% Formula description
		\node[below=0.2cm of Result] {\small $\pi_{R \cup S}(\sigma_{R.A = S.A}(R \times S))$};
		
	\end{tikzpicture}
	\caption{Pictorial representation of Relational Algebra Natural Join}
	\label{fig:natural_join}
\end{figure}
\begin{align}
	R \bowtie S = \pi_{R \cup S}\brak{\sigma_{R.A = S.A}\brak{R \times S}} \label{eq:DA2025-join}
\end{align}
\solution
The equation describes the mechanical steps of combining two database tables using relational algebra operators:
\begin{itemize}
	\item $R$ and $S$: The two input relations (data tables) being operated on.
	\item $\bowtie$ (Natural Join): The overarching operator that merges $R$ and $S$ by matching rows with identical values in their common columns, outputting a single combined table with no duplicate columns.
	\item $R \times S$ (Cartesian Product): The foundational operation that pairs every single row in table $R$ with every single row in table $S$, generating all possible combinations.
	\item $\sigma_{R.A = S.A}$ (Selection): The row-filtering step. The $\sigma$ operator evaluates the Cartesian product and retains only the rows where the value in the shared attribute $A$ matches between both tables.
	\item $\pi_{R \cup S}$ (Projection): The column-filtering step. The $\pi$ operator dictates which columns appear in the final result. The subscript $R \cup S$ ensures that all unique columns from both tables are kept, effectively stripping out the redundant duplicate common column ($S.A$) generated by the Cartesian product.
\end{itemize}
% =========================================================
% Q.18
% [Reference: General Domain Knowledge - Data Structures (Hashing)]
% =========================================================
\item Linear Probing in Hash Tables
% linear_probing.tex
\begin{figure}[H]
	\centering
	\begin{tikzpicture}[
		slot/.style={draw, minimum width=3.5cm, minimum height=0.8cm, align=center},
		occupied/.style={slot, fill=gray!30},
		empty/.style={slot, fill=white},
		>=Stealth,
		start chain=1 going right,
		node distance=-\pgflinewidth
		]
		
		% Array structure
		\node [on chain=1] {\dots};
		\node [on chain=1, occupied] (slot0) {Occupied};
		\node [on chain=1, occupied] (slot1) {Occupied};
		\node [on chain=1, empty]    (slot2) {Empty};
		\node [on chain=1] {\dots};
		
		% Indices below the array
		\node[below=0.3cm of slot0] {$h(k) \bmod M$};
		\node[below=0.3cm of slot1] {$h(k)+1 \bmod M$};
		\node[below=0.3cm of slot2] {$h(k)+2 \bmod M$};
		
		% Key input
		\node[above=2.5cm of slot0, draw, rounded corners, inner sep=0.2cm] (key) {Insert Key $k$};
		
		% Probing Arrows
		\draw[->, thick] (key.south) -- (slot0.north)
		node[midway, left=0.1cm] {$i=0$};
		
		\draw[->, thick] (slot0.north) .. controls +(up:1.5cm) and +(up:1.5cm) .. (slot1.north)
		node[midway, above] {$i=1$ (Collision)};
		
		\draw[->, thick] (slot1.north) .. controls +(up:1.5cm) and +(up:1.5cm) .. (slot2.north)
		node[midway, above] {$i=2$ (Success)};
		
	\end{tikzpicture}
	\caption{Pictorial representation of Linear Probing in Hash Tables}
	\label{fig:linear_probing}
\end{figure}
\begin{align}
	h\brak{k, i} = \brak{h\brak{k} + i} \pmod M \label{eq:DA2025-probing}
\end{align}
The terms in the probing function represent:
\begin{itemize}
	\item $k$: The original key or piece of data to be stored.
	\item $h\brak{k}$: The primary hash function that calculates the initial target index.
	\item $i$: The probe attempt number, starting at $0$ and incrementing linearly for each collision encountered.
	\item $M$: The total allocated size of the hash table.
	\item $\pmod M$: The modulo operation ensuring the probe sequence wraps around to index $0$ if the end of the table is reached.
	\item $h\brak{k, i}$: The final calculated storage index for the $i$-th probe attempt.
\end{itemize}
% =========================================================
% Q.19
% [Reference: formulae.tex (Item 8: Median definition adjusted for CDFs)]
% =========================================================
\item Median of a Continuous Distribution

Let $f_X\brak{x}$ be the probability density function. The median $m$ divides the total probability mass strictly in half:
\begin{align}
	\pr{X \le m} &= 0.5
	\\
	\int_{-\infty}^{m} f_X\brak{x} \, dx &= 0.5
\end{align}
Because the integral of $f_X\brak{x}$ from $-\infty$ to $m$ is exactly $F_X\brak{m}$, it follows that:
\begin{align}
	F_X\brak{m} = \int_{-\infty}^{m} f_X\brak{x} \, dx = 0.5 \label{eq:DA2025-median}
\end{align}


% =========================================================
% Q.20
% [Reference: formulae.tex (Item 16: Gaussian Functions & Variances)]
% =========================================================
\item Let the random variable $Z$ follow a standard normal distribution.
\begin{align}
	\intertext{It follows that}
	E\sbrak{Z} &= 0 \label{eq:DA2025-E_Z}
	\\
	E\sbrak{Z^2} &= 1 \label{eq:DA2025-E_Z2}
	\intertext{Show that for}
	X &= aZ + b \label{eq:DA2025-X_def}
	\intertext{It follows that}
	E\sbrak{X} &= aE\sbrak{Z} + b
	\\
	\text{Var}\brak{X} &= a^2\text{Var}\brak{Z} \label{eq:DA2025-affine_rv}
\end{align}

\solution
The step-by-step derivations for the expectations are:
\begin{align}
	E\sbrak{X} &= E\sbrak{aZ + b}
	\\
	&= aE\sbrak{Z} + b
	\\
	&= a\brak{0} + b \quad \text{(using \eqref{eq:DA2025-E_Z})}
	\\
	&= b
	\intertext{Then,}
	X - E\sbrak{X} &= \brak{aZ + b} - b
	\\
	&= aZ
	\intertext{Similarly,}
	E\sbrak{\brak{X - E\sbrak{X}}Z} &= E\sbrak{aZ \cdot Z}
	\\
	&= aE\sbrak{Z^2}
	\\
	&= a\brak{1} \quad \text{(using \eqref{eq:DA2025-E_Z2})}
	\\
	&= a
	\intertext{Yielding,}
	E\sbrak{\brak{X - E\sbrak{X}}^2} &= E\sbrak{\brak{aZ}^2}
	\\
	&= a^2E\sbrak{Z^2}
	\\
	&= a^2\brak{1} \quad \text{(using \eqref{eq:DA2025-E_Z2})}
	\\
	&= a^2
\end{align}
	% =========================================================
	% Q.21
	% [Reference: bt-2025-formulae.tex (Item 3: Integrals for Laplace logic)]
	% =========================================================
	\item Tail Probability of Exponential Distribution
	\begin{align}
		\pr{X \ge x} = \int_{x}^{\infty} \lambda e^{-\lambda t} \, dt = e^{-\lambda x} \label{eq:DA2025-exp_tail}
	\end{align}
	
	% =========================================================
	% Q.23, Q.47
	% [Reference: General Domain Knowledge - Python sets and lists]
	% =========================================================
	% =========================================================
	% Q.23
	% [Reference: General Domain Knowledge - Python Data Structures]
	% =========================================================
	\item Python List and Set Operations
	\begin{align}
		A.\text{update}\brak{B} \implies A \gets A \cup B \label{eq:DA2025-set_update}
	\end{align}
	\solution
	Let $A = \brak{a_1, a_2, \dots, a_n}$ and $B = \brak{b_1, b_2, \dots, b_m}$ represent sequences, and let $\gets$ denote in-place modification. The functions map to the following mathematical transformations[cite: 41]:
	\begin{align}
		A.\text{extend}\brak{B} &\implies A \gets \brak{a_1, \dots, a_n, b_1, \dots, b_m}
		\\
		A.\text{append}\brak{B} &\implies A \gets \brak{a_1, \dots, a_n, B}
		\\
		A.\text{update}\brak{B} &\implies A \gets A \cup B \quad \text{(where $A$ and $B$ are sets)}
		\\
		A.\text{insert}\brak{i, B} &\implies A \gets \brak{a_1, \dots, a_i, B, a_{i+1}, \dots, a_n}
	\end{align}
	% =========================================================
	% Q.24
	% [Reference: bt-2025-formulae.tex (Item 10: Continuity and Differentiability)]
	% =========================================================
	\item Differentiability Rules
	\begin{align}
		\brak{fg}'\brak{c} = f'\brak{c}g\brak{c} + f\brak{c}g'\brak{c} \label{eq:DA2025-product_rule}
	\end{align}
	\solution
	If $f$ and $g$ are differentiable at $c$, their sum, difference, and product are strictly differentiable at $c$. The quotient is differentiable if $g\brak{c} \ne 0$. Composition $f \circ g$ requires $f$ to be differentiable at $g\brak{c}$, which is not universally guaranteed.
	
	% =========================================================
	% Q.25
	% [Reference: ae-2026_formulae.tex (Item 16: 2-D Basis Vectors)]
	% =========================================================
	\item Properties of Orthonormal Sets
	\begin{align}
		\vec{u}_i^\top \vec{u}_j = \delta_{ij} \label{eq:DA2025-orthonormal}
	\end{align}
	\solution
	An orthonormal set of vectors in $\mathbb{R}^n$ is unconditionally linearly independent. Orthonormal bases in $\mathbb{R}^n$ are not unique because applying any rotation matrix generates a new valid orthonormal basis.
	
	% =========================================================
	% Q.27
	% [Reference: General Domain Knowledge - Algorithms/Searching]
	% =========================================================
	\item Binary Search Complexity
	\begin{align}
		T\brak{n} = T\brak{n/2} + O\brak{1} \implies T\brak{n} = O\brak{\log n} \label{eq:DA2025-binary_search}
	\end{align}
	\solution
	Binary search requires $O\brak{1}$ random access to indices, restricting its $O\brak{\log n}$ performance bounds strictly to contiguous arrays, not linked lists.
	
	% =========================================================
	% Q.28
	% [Reference: ae-2026_formulae.tex (Item 20: Matrix Determinant Lemma)]
	% =========================================================
	\item Matrix Determinant Lemma
	\begin{align}
		\mydet{\vec{I} + \vec{u}\vec{v}^\top} = 1 + \vec{v}^\top\vec{u} \label{eq:DA2025-det_lemma}
	\end{align}
	\solution
	For $\vec{A} = \vec{I}_n + \vec{x}\vec{x}^\top$ with $\vec{x}^\top\vec{x} = 1$:
	\begin{align}
		\mydet{\vec{A}} &= 1 + \vec{x}^\top\vec{x} = 2 \ne 0
	\end{align}
	Because $\mydet{\vec{A}} \ne 0$, $\vec{A}$ is invertible and has full rank. Its eigenvalues are $2$ and $1$, neither being zero or negative.
	
	% =========================================================
	% Q.29
	% [Reference: General Domain Knowledge - Sorting algorithms]
	% =========================================================
	\item Insertion Sort Inversions
	\begin{align}
		\text{Swaps} = \text{Number of Inversions} \label{eq:DA2025-inversions}
	\end{align}
	\solution
	Insertion sort executes one swap per inverted pair. For the array to be sorted in exactly two swaps, the inserted element $x$ must be strictly smaller than exactly two preceding elements.
	
	% =========================================================
	% Q.31
	% [Reference: ce2-2025-problems.tex (Q.22: Conditional Probability Rules)]
	% =========================================================
	\item Bayes' Theorem for Partitioned Spaces
	\begin{align}
		\pr{B_j \mid W} = \frac{\pr{W \mid B_j}\pr{B_j}}{\sum_i \pr{W \mid B_i}\pr{B_i}} \label{eq:DA2025-bayes}
	\end{align}
	
	% =========================================================
	% Q.32
	% [Reference: formulae.tex (Item 7: Definition of Derivative / Limits)]
	% =========================================================
	\item Asymptotic Radical Limits
	\begin{align}
		\lim_{t \to \infty} \brak{\sqrt{t^2 + at} - t} = \frac{a}{2} \label{eq:DA2025-radical_limit}
	\end{align}
	\solution
	Multiplying by the conjugate:
	\begin{align}
		\lim_{t \to +\infty} \frac{\brak{t^2 + t} - t^2}{\sqrt{t^2 + t} + t} &= \lim_{t \to +\infty} \frac{1}{\sqrt{1 + \frac{1}{t}} + 1} = 0.5
	\end{align}
	
	% =========================================================
	% Q.34
	% [Reference: bt-2025-formulae.tex (Item 6: First-Order Parameter Update/Gradient)]
	% =========================================================
	\item Least Squares Regression Through Origin
	\begin{align}
		w^* = \frac{\sum x_i y_i}{\sum x_i^2} \label{eq:DA2025-least_squares}
	\end{align}
	\solution
	Minimizing $S\brak{w} = \sum \brak{y_i - w x_i}^2$:
	\begin{align}
		\frac{dS}{dw} &= -2\sum x_i \brak{y_i - w x_i} = 0 \implies w^* = \frac{24}{14} \approx 1.714
	\end{align}
	
	% =========================================================
	% Q.35
	% [Reference: ce2-2025-problems.tex (Q.22: Conditional Probability Rules)]
	% =========================================================
	\item Bayes Optimal Error Rate
	\begin{align}
		\pr{\text{error} \mid \vec{x}} = \frac{\min_k \pr{\vec{x} \mid y_k}\pr{y_k}}{\pr{\vec{x}}} \label{eq:DA2025-bayes_error}
	\end{align}
	\solution
	The classifier selects the class with the maximum posterior probability. The error probability is determined by the unselected (minimum) posterior:
	\begin{align}
		\min\cbrak{\frac{3}{4} \cdot \frac{1}{3}, \frac{1}{4} \cdot \frac{2}{3}} &= \frac{1}{6}
		\\
		\pr{\vec{x}} &= \frac{1}{4} + \frac{1}{6} = \frac{5}{12}
		\\
		\pr{\text{error} \mid \vec{x}} &= \frac{1/6}{5/12} = 0.40
	\end{align}
	
	% =========================================================
	% Q.36
	% [Reference: formulae.tex (Item 16: Gaussian properties)]
	% =========================================================
	\item Moments of Standard Normal and Chi-Square
	\begin{align}
		Z \sim \mathcal{N}\brak{0,1} \implies E\sbrak{Z^2} = 1, \; E\sbrak{Z^4} = 3 \label{eq:DA2025-chisq}
	\end{align}
	\solution
	For $Y = Z^2$:
	\begin{align}
		\text{Var}\brak{Y} &= E\sbrak{Y^2} - \brak{E\sbrak{Y}}^2 = E\sbrak{Z^4} - \brak{E\sbrak{Z^2}}^2 = 3 - 1^2 = 2
	\end{align}
	
	% =========================================================
	% Q.37
	% [Reference: ae-2026_formulae.tex (Item 18: Matrix properties)]
	% =========================================================
	\item Rank of Involutory Matrices
	\begin{align}
		\vec{A}^3 = \vec{A} \implies \text{rank}\brak{\vec{A}} = \text{rank}\brak{\vec{A}^2} \label{eq:DA2025-annihilating}
	\end{align}
	\solution
	Factoring $\vec{A}^3 - \vec{A} = \vec{A}\brak{\vec{A} - \vec{I}}\brak{\vec{A} + \vec{I}} = \vec{0}$. The minimal polynomial has distinct linear roots, making $\vec{A}$ diagonalizable. The Rank-Nullity theorem confirms equal rank across powers.
	
	% =========================================================
	% Q.38
	% [Reference: bt-2025-formulae.tex (Item 8: System of Linear Equations)]
	% =========================================================
	\item Gram Matrix Positive Definiteness
	\begin{align}
		\vec{A} = \vec{X}^\top\vec{X} \implies \vec{z}^\top\vec{A}\vec{z} = \norm{\vec{X}\vec{z}}^2 \label{eq:DA2025-gram}
	\end{align}
	\solution
	If the columns of $\vec{X}$ are linearly independent, $\vec{X}\vec{z} = \vec{0} \iff \vec{z} = \vec{0}$. Thus, $\norm{\vec{X}\vec{z}}^2 > 0$ for all non-zero $\vec{z}$, rendering $\vec{A}$ positive definite and invertible.
	
	% =========================================================
	% Q.39
	% [Reference: General Domain Knowledge - Probability CDF mapping]
	% =========================================================
	\item Interval Probabilities from CDF
	\begin{align}
		\pr{X^2 \le c^2} = F_X\brak{c} - F_X\brak{-c} \label{eq:DA2025-interval_cdf}
	\end{align}
	\solution
	For $F_X\brak{x} = \frac{1}{4}\brak{x + 1}^2$:
	\begin{align}
		\pr{X^2 \le 0.25} &= F_X\brak{0.5} - F_X\brak{-0.5} = \frac{2.25}{4} - \frac{0.25}{4} = 0.5
	\end{align}
	
	% =========================================================
	% Q.40
	% [Reference: formulae.tex (Item 11: Binomial Distribution & Item 17: CLT)]
	% =========================================================
	\item Central Limit Theorem for Binomial Sums
	\begin{align}
		Z_n = \frac{\sum X_i - n\theta}{\sqrt{n\theta\brak{1-\theta}}} \xrightarrow{d} \mathcal{N}\brak{0,1} \label{eq:DA2025-clt}
	\end{align}
	\solution
	For $n = 300, \theta = 0.25$, $\mu = 75$ and $\sigma = 7.5$:
	\begin{align}
		\pr{60 \le Y \le 90} &\approx \phi\brak{\frac{90-75}{7.5}} - \phi\brak{\frac{60-75}{7.5}} = \phi\brak{2} - \phi\brak{-2}
	\end{align}
	
	% =========================================================
	% Q.41
	% [Reference: bt-2025-formulae.tex (Item 3: Integrals)]
	% =========================================================
	\item Geometric Discretization of Exponential
	\begin{align}
		\pr{\lfloor X \rfloor = l} = e^{-\lambda l}\brak{1 - e^{-\lambda}} \label{eq:DA2025-floor_exp}
	\end{align}
	\solution
	Integrating $\lambda e^{-\lambda x}$ over $[l, l+1]$ yields $q^l\brak{1-q}$ with $q = e^{-\lambda}$.
	
	% =========================================================
	% Q.45
	% [Reference: ce2-2025-problems.tex (Q.22: Probability basics)]
	% =========================================================
	\item Complement Rule
	\begin{align}
		\pr{\text{at least one success}} = 1 - \brak{1 - p}^n \label{eq:DA2025-complement}
	\end{align}
	\solution
	The probability of rolling zero $1$s on 100 fair dice is $\brak{5/6}^{100}$.
	
	% =========================================================
	% Q.49
	% [Reference: ce2-2025-problems.tex (Q.38: Local Extrema)]
	% =========================================================
	\item Critical Points
	\begin{align}
		f'\brak{x} = 0 \text{ at local extrema} \label{eq:DA2025-extrema}
	\end{align}
	\solution
	Local extrema require evaluating roots of $f'\brak{x}$ and checking $f''\brak{x}$ convexity. Extrema of the derivative itself locate inflection points.
	
	% =========================================================
	% Q.50
	% [Reference: ae-2026_formulae.tex (Item 16: Orthogonal sets)]
	% =========================================================
	\item Orthogonal Projection Matrices
	\begin{align}
		\vec{A} = \sum_{i=1}^k \vec{x}_i \vec{x}_i^\top \implies \vec{A}^2 = \vec{A} \label{eq:DA2025-projection}
	\end{align}
	\solution
	Using orthonormal properties, $\vec{A}$ is idempotent. Singular values equal eigenvalues $\sigma \in \{0, 1\}$. For incomplete bases $k < n$, $\vec{A}$ is singular.
	
	% =========================================================
	% Q.51
	% [Reference: bt-2025-formulae.tex (Item 6: Convexity)]
	% =========================================================
	\item Strict Convexity
	\begin{align}
		f''\brak{x} > 0 \implies \text{Unique global minimum (if any)} \label{eq:DA2025-convex}
	\end{align}
	\solution
	A strictly monotonic $f'\brak{x}$ cannot possess multiple roots, precluding multiple local minima.
	
	% =========================================================
	% Q.52
	% [Reference: ae-2026_formulae.tex (Item 18: Matrix Determinant/Norm)]
	% =========================================================
	\item Isometric Norm Preservation
	\begin{align}
		\norm{\vec{Ax}}^2 = \norm{\vec{x}}^2 \iff \vec{A}^\top\vec{A} = \vec{I} \label{eq:DA2025-isometry}
	\end{align}
	\solution
	Applying polarization identity, the equality strictly implies orthogonality and full rank.
	
	% =========================================================
	% Q.53
	% [Reference: General Domain Knowledge - SVM Optimization]
	% =========================================================
	\item SVM Margin
	\begin{align}
		\gamma = \frac{2}{\norm{\vec{w}}} \label{eq:DA2025-svm}
	\end{align}
	\solution
	Support vectors lie directly on the canonical hyperplanes $y_i\brak{\vec{w}^\top\vec{x}_i + b} = 1$. The geometric width is inverse to norm $\vec{w}$.
	
	% =========================================================
	% Q.54
	% [Reference: formulae.tex (Item 17: Central Limit Theorem / Sample Means)]
	% =========================================================
	\item Variance of Sample Mean
	\begin{align}
		\text{Var}\brak{\frac{1}{n}\sum_{i=1}^n X_i} = \frac{\sigma^2}{n} \label{eq:DA2025-sample_mean}
	\end{align}
	\solution
	Variance of a sum of independent identical variables scales by $n$, rendering the sample mean variance inversely proportional to $n$.
	
	% =========================================================
	% Q.58
	% [Reference: ce2-2025-problems.tex (Q.53: Network/Graph Constraints)]
	% =========================================================
	\item Shortest Path Metric
	\begin{align}
		d\brak{u,v} \le d\brak{u,w} + d\brak{w,v} \label{eq:DA2025-triangle_ineq}
	\end{align}
	\solution
	Fake edges violating the triangle inequality are identifiable by calculating the minimum spanning distances directly.
	
	% =========================================================
	% Q.59
	% [Reference: formulae.tex (Item 7: First Principles Derivative)]
	% =========================================================
	\item Lipschitz Continuity with Exponent $> 1$
	\begin{align}
		\abs{f\brak{x} - f\brak{y}} \le K\abs{x - y}^\alpha \implies f'\brak{x} = 0 \label{eq:DA2025-lipschitz}
	\end{align}
	\solution
	Dividing by $\abs{x-y}$ and taking the limit $y \to x$ for $\alpha > 1$ yields a derivative of zero, enforcing a constant function.
	
	% =========================================================
	% Q.60
	% [Reference: ae-2026_formulae.tex (Item 18: Eigenvalues trace mapping)]
	% =========================================================
	\item Rayleigh-Ritz Theorem for PCA
	\begin{align}
		\max_{\norm{\vec{u}}=1} \vec{u}^\top\vec{S}\vec{u} = \lambda_{\max}\brak{\vec{S}} \label{eq:DA2025-pca}
	\end{align}
	\solution
	Empirical variance along $\vec{u}$ is maximized by the eigenvector corresponding to the principal eigenvalue.
	
	% =========================================================
	% Q.61
	% [Reference: bt-2025-formulae.tex (Item 9: Sum of Discrete RVs)]
	% =========================================================
	\item Expectation of Indicator Sums
	\begin{align}
		E\sbrak{\sum_{i=1}^n I_i} = \sum_{i=1}^n E\sbrak{I_i} = np \label{eq:DA2025-linearity}
	\end{align}
	\solution
	Linearity of expectation holds uniformly irrespective of variable dependence or independence.
	
	% =========================================================
	% Q.62
	% [Reference: General Domain Knowledge - SQL/Relational Algebra]
	% =========================================================
	\item Relational Division
	\begin{align}
		R \div S = \pi_A\brak{R} - \pi_A\brak{\brak{\pi_A\brak{R} \times S} - R} \label{eq:DA2025-rel_div}
	\end{align}
	\solution
	Simulates universal quantification $\forall$ in relational schema.
	
\end{enumerate}